\documentclass[12pt,a4paper]{article}
\usepackage[utf8]{inputenc}
\usepackage{amsmath}
\usepackage{amsfonts}
\usepackage{amssymb}
\usepackage{geometry}
\usepackage{booktabs}

\geometry{margin=1in}

\begin{document}

\begin{center}
    \large \textbf{The Mechanical Derivation of Light Speed: \\ Phase Velocity as a Resultant of Vacuum Elasticity} \\
    \vspace{1em}
    \normalsize Robin Skavberg \\
    \normalsize February 2026
\end{center}

\begin{abstract}
This paper derives the speed of light (c) not as an axiomatic constant, but as a mechanical property of the vacuum medium. By defining the vacuum as an elastic manifold with an intrinsic restorative stiffness, we demonstrate that the propagation velocity of electromagnetic waves is governed by the local matter-energy load ($\rho_m$). This derivation provides a unified framework for understanding the 1.7s multi-messenger delay and suggests that the "universal speed limit" is a manifestation of the vacuum's ground-state impedance.
\end{abstract}

\section{The Vacuum as an Elastic Medium}
In classical wave mechanics, the velocity of a wave ($v$) through a medium is defined by the ratio of the medium's stiffness ($k$) to its inertia or density ($\rho$):
\begin{equation}
v = \sqrt{\frac{k}{\rho}}
\end{equation}
Applying this to the vacuum, we define the restorative stiffness as $c^2$ and the gravitational coupling $G$ as the facilitator of field response.

\section{The Derivation of Variable Velocity}
We begin by defining the relationship between the gravitational constant and the vacuum load:
\begin{equation}
G = \frac{c^2}{8\pi\rho_m}
\end{equation}
This implies that matter-energy density ($\rho_m$) acts as the inertial "load" on the vacuum's elastic response.

\subsection{Phase Velocity and Refractive Index}
The speed of light in a vacuum ($v_{\gamma}$) is modified by the Presence of matter, creating a refractive index $n_v$:
\begin{equation}
n_v = \sqrt{1 + \rho_m}
\end{equation}
Substituting the definition of $\rho_m$ from Eq. (2) into the index:
\begin{equation}
n_v = \sqrt{1 + \frac{c^2}{8\pi G}}
\end{equation}

\subsection{The Final Velocity Equivalence}
The propagation velocity is given by $v_{\gamma} = c / n_v$. Substituting Eq. (4) yields:
\begin{equation}
v_{\gamma} = \frac{c}{\sqrt{1 + \frac{c^2}{8\pi G}}}
\end{equation}
Reducing this expression to its simplest form:
\begin{equation}
v_{\gamma} = c \sqrt{\frac{8\pi G}{8\pi G + c^2}}
\end{equation}



\section{Physical Significance}
\begin{itemize}
    \item \textbf{Ground-State Limit:} In a pure vacuum ($\rho_m \to 0$), $G$ reaches its theoretical maximum, and $v_{\gamma}$ approaches $c$. This defines $c$ as the "stiffness limit" of the universe.
    \item \textbf{Vacuum Refraction:} Curvature is reinterpreted as a gradient in the speed of light. Light does not bend because space is "curved" in a geometric sense, but because the vacuum impedance changes, refracting the wave path.
    \item \textbf{Multi-Messenger Symmetry:} Gravitational waves, being oscillations of the stiffness modulus itself, propagate at $c$. Light, being an electromagnetic excitation within the loaded medium, propagates at $v_{\gamma}$. The 1.7s delay (GW170817) is the time-integral of this velocity discrepancy.
\end{itemize}



\section{Conclusion}
The speed of light is a dynamic variable in equilibrium with the tensor strength of the vacuum. Equation (6) provides the mathematical link between the gravitational constant and electromagnetic propagation, removing the need for Dark Matter by accounting for the vacuum's restorative response in low-density environments.

\begin{thebibliography}{9}
\bibitem{1} Skavberg, R. (2026). \textit{The Mechanical Prediction for Galactic Rotation Curves}. 
\bibitem{2} Maxwell, J. C. (1865). \textit{A Dynamical Theory of the Electromagnetic Field}.
\bibitem{3} Abbott, B. P., et al. (2017). \textit{Multi-messenger Observations of a Binary Neutron Star Merger}.
\end{thebibliography}

\end{document}